\documentclass[11pt]{article}
\usepackage{../shared/luxcover}
\usepackage{amsmath,amssymb}
\usepackage{hyperref}
\usepackage{booktabs}
\usepackage{enumitem}
\usepackage{xcolor}

\title{Lux M-Chain: Regulated Real-World Assets}
\author{Zach Kelling\\\textit{Lux Industries, Inc.}\\\texttt{research@lux.network}}
\date{December 2023}

\begin{document}

\luxcoverpage

\begin{abstract}
Regulated Blockchain Layer for the Creation Of RWA Tokens

Bringing the power of smart contracts to banks and regulated financial institutions, MChain serves to remove friction for companies and applications that seek the automation of regulatory requirements, increasing their speed and efficiency in the creation of securities and other types of financial instruments.
\end{abstract}

\tableofcontents
\newpage


Bringing the power of smart contracts to banks and regulated financial institutions, MChain serves to remove friction for companies and applications that seek the automation of regulatory requirements, increasing their speed and efficiency in the creation of securities and other types of financial instruments.


Lux MChain aims for widespread acceptance in both business and society through the automation of trustless interactions and has created an efficient means to organize the creation and authentication of real world asset tokens.


Lux MChain has partnered with CDAX to provide an integrated money transmitters license to address the needs of the financial services sector while also solving the critical issues of scalability, privacy, and interoperability that stand in the way of wide-scale mass adoption.


\subsubsection{Bridging The Gap to Put The People In Power Over Money Creation}

MChain bridges the gap for a user base and society that is used to having their value secured by a third party while at the same time adhering to regulation from the judiciary oversight mechanisms put in place to deter and prevent fraud and abuse of the network. Oversight mechanisms also promote compliance with ethical, statutory, and regulatory standards.


DAO governance and community voting for RWA tokens allows for the owners of those assets to choose how they are used, staked or sold. Ownership of the RWA gives the community the power to use the technology to create a transparent and fair ecosystem of money creation, by allowing them to make decisions on how their asset is used and how it is managed. This allows the people to have direct control and influence over the asset and the money supply created from the asset. It provides the community the ability to hold others accountable for the decisions they make in regards to the money created on the MChain platform. This helps ensure a secure, transparent and fair environment.


\subsubsection{Benefits of MChain}

Provides banking and financial institutions with a framework to ensure compliance with various local and international regulations.


Allows companies to gain access to banking and financial services that are required to meet regulatory requirements.


Ensures the efficient and secure transfer of funds between different countries and jurisdictions.


Eliminates the risk of interacting with unregulated financial products.


Automated KYC/AML to ensure that all participants are accredited investors


\section{Problem with Current Blockchain Solutions}

The current blockchain environment, while revolutionary in many ways, has not yet provided a solution that is both quantum resistant and scalable to billions of users, while providing true privacy for transactions. As a result, there is a need for new blockchain solutions that can meet these needs and provide a more secure, private, and efficient environment.


\subsection{Lux MChain address these three main concerns with the following solutions:}

SECURE --- Lux wallets leverage post quantum cryptography to protect private keys


PRIVATE --- Zero knowledge bridge creates a hashed reduction of all transactions


SCALABLE --- Lux Network Subnet Architecture allows Billions of Users to Scale >1m/TPS


COMPLIANT --- KYC/AML procedures implemented


\subsection{Regulatory Compliance and Licensing}

An important part in creating such solutions is understanding that compliance and safety have been secured through trusted institutions and regulatory bodies answerable to governments and people. To bridge the gap between the current reality and the uncertain decentralized future, it is essential to embrace the role that regulation has played and ensure the same security and safety that users have come to expect is present in the new technology. It's the transparency of the audits and the viability into the creation of money that will set us free from the slavery of debt.


Prior to blockchain, these rights have been secured by regulatory bodies. It must be understood that the same security that people have come to expect, is the same environment this new technology must exist under. It is also the responsibility and obligation of the creators of this new financial ecosystem to provide the trustless and decentralized infrastructure that promotes true democracy and true freedom of the people, as it restores the control mechanisms into the hands of the people that are the holders to the burden of public debt. No matter what the new solution is, it must be understood that the control mechanisms must remove the control of central authority.


It is this expectation of a regulated environment that has made the public susceptible to abuse, fraud and hacks.


\section{A Multifaceted problem}

The problem has two sides; on one hand, a public that is used to trust and security being assured by a centralized power, and on the second hand, a lack of compliant, transparent and regulated products to offer.


MChain meets both needs; by providing decentralized products that are compliant and working with the current regulatory bodies and that are safe enough to satisfy consumers and government expectations. This requires consistent transparency to physical audits and automated security requirements to ensure that the underlying value remains safe.


\section{The Oracle Problem}

The current blockchain ecosystem depends on a network of unaffiliated exchanges that allow people to transfer between various blockchains and cryptocurrencies, as well as into and out of various real world assets such as physical assets and currencies. There is a problem with these centralized exchanges including high fees, opaque order books, and questionable criteria for listing on exchanges (i.e., who gets listed). The biggest hurdle to current blockchain ecosystems is the reliability of off chain data being validated before being brought on chain.


Many of blockchain's innovations and issues are informed and distorted by existing real world solutions. Exchanges are an example of that. We're accustomed to trusting a centralized exchange for the conversion of fungible things. We're accustomed to paying a fee and outsourcing the mechanism of that transaction to a third party (i.e., a bank).


Blockchain and its inherent trust and lowest cost nature was designed to obviate the need for these centralized authorities. Lux executes upon this function at its most fundamental and at the lowest cost possible, and seeks to provide a decentralized, yet regulatory compliant path forward.


\subsection{Exchange Problem}

Exchanges are a major weak point in the blockchain sphere, and present a significant risk to users. Issues such as front running, wallet theft, bridge vulnerabilities, compromised identities, and opaque fees and costs are just a few of the challenges that threaten to erode the promises that blockchain technology offers such as decentralization and cost reduction. To meet these goals, new solutions that enable greater security and transparency are essential.


\section{Gateway Between the Physical and Digital}

The relationship between cryptocurrency and financial ledgers is important yet tenuous. For the most part, all real world costs are settled in fiat currencies. Any friction, complication, or delay in cryptocurrency adoption and use is seen as fatal by consumers.


Mchain provides an easier and more cost effective way to move between physical assets and blockchain-based digital assets. It does this by providing a platform that is blockchain agnostic with built-in compliance features such as Know Your Customer (KYC) and Anti Money Laundering (AML) requirements. These features address some of the limitations that prevent blockchain technology from being adopted more widely, and make it more consistent with the traditional financial system.


\subsubsection{Key Features of MChain's Role in SWIFT Integration for Asset Verification}

Lux Exchange's integration with SWIFT messaging, facilitated by MChain, plays a crucial role in verifying assets, particularly in transactions involving vaulted assets. This feature is especially significant for institutional participants and investors engaged in high-value transactions.


\subsection{Utilizing MT-760 and Different Types of SWIFT Codes}

\subsection{MT-760 for Asset Verification:}

Role in Issuing Letters of Credit: MT-760 is a SWIFT message format used primarily by banks for issuing or requesting letters of credit or documentary credits​​​​.


Blocking Funds for Transactions: When an MT-760 message is sent, it indicates that the issuing institution will ring-fence the necessary funds for the transaction, essentially serving as a guarantee of payment​​.


MT-799 for Pre-Advice and Proof of Funds:


Pre-advice Communication: MT-799 acts as a pre-advice communication tool between banks, often serving as a confirmation message or proof of funds​​.


Non-Transactional Nature: Unlike MT-760, MT-799 is not used for transferring funds but for showing proof of deposits without any commitment to transfer funds​​.


\subsection{Advantages of SWIFT Integration via MChain}

Enhanced Verification and Trust: The use of MT-760 and MT-799 SWIFT codes ensures a higher level of verification, enhancing trust and transparency in transactions, particularly important for institutional and high-net-worth participants.


Seamless Fiat Onboarding and Offboarding: MChain facilitates the conversion of fiat currencies into cryptocurrencies and vice versa, broadening Lux Exchange's accessibility to a wider array of users, including those familiar with traditional banking systems.


Regulatory Compliance: Adherence to international banking regulations is ensured, making transactions compliant and secure, essential for maintaining the platform's integrity and user trust.


Streamlined Institutional Participation: The integration is advantageous for institutional investors, providing a familiar and secure banking protocol for engaging with the crypto market, thus bridging traditional finance with decentralized finance (DeFi).


Innovative Financial Bridge: MChain's role in integrating SWIFT messages is a testament to Lux Exchange's commitment to innovation, melding the rapidly evolving DeFi space with established financial systems.


\section{Technology Innovations}

\subsection{Fully decentralized, non-custodial platform}

Digital Assets are never in our possession.


You are in full control of your private keys at all times.


Federated trust lines between various regulated entities


\subsection{Blockchain agnostic}

Zero Knowledge, Quantum Teleport cross-chain swaps


Support for multiple chains and liquidity pools


\subsection{Compliance Innovations}

Automated Compliance Notifications


Physical Asset Redemptions


Automated  KYC/AML Registry


Regulatory Transparency


\subsection{Versioned and upgradeable protocols}

Easily adapt to changing compliance / regulatory needs


Quickly adopt rapidly developing protocol standards


\subsection{Scaling via ZK Proofs}

Execute smart contracts off-chain, blockchain enforces smart contract code only in case of disagreement


Improved privacy, scalability


\subsection{Oracles}

Lux Mchain leverages Lux Oracle to provide reliable off-chain data. Integrate pre-built, time-tested oracle solutions that has already secured smart contracts for market-leading decentralized applications.


Regulatory Compliance integrated Into MChain Security


Connect real world data and third-party systems with smart contract code


Automated interest disbursements


Automated Physical Asset Audit for Validation of RWA Tokens


Authentication of Geological Surveys for creation of RWA Tokens


\section{Near Zero Fee Products}

By lowering fees at the protocol layer and allowing liquid, compliant flexible ownership of the underlying asset, these products are finally able to deliver on one of blockchain's ultimate promises: the reduction of fees through automation and the elimination of the cost of centralizing trust. To date, the solutions have been orders of magnitudes more expensive than the traditional financial service products. Lux allows for the deployment of the structurally lowest fee structure of financial service products possible and finally realizing the ambition of efficient low cost blockchain products.


\section{The Connective Tissue of the Blockchain Space}

Until now, the blockchain space has been a collection of competing and unrelated protocol. The vision of fast, cheap decentralized solutions to our current centralized world has not materialized. While the current financial system is centralized, it is also fast and reliable. This is accomplished by agreed upon communication and security protocols like SWIFT. How do we achieve such harmony and communication between the ever evolving and independant areas of the blockchain ecosystem? Lux offers a potential solution to this basic and critical functionality. By offering this connective piece at the lowest possible cost in both time and speed, while being protocol andgostic, Lux has provided the universal underpinning that is necessary for this concept to reach the diversity scale that is possible.


\section{Lux Treasury Token}

\subsection{Underlying Structure}

Legally, the Lux Treasury Token will be digitized shares of a registered investment company formed pursuant to the Investment Company Act of 1940.  However, unlike traditional mutual funds, the Lux Treasury Token will be closed-end, with regular and frequent redemption periods.  We anticipate the legal structure will be sufficiently robust to support peer-to-peer transactions as well as transactions on exchanges.  Since the underlying assets will be actual US Treasury securities, which are very liquid, the fund will be able to meet any redemption requests to exchange the shares for fiat.  All such requests will be processed through a registered custodian (such as a U.S. bank) and wired directly to the bank account of the investor.


\subsection{Fully Compliant and Regulated}

The legal structure will be fully compliant with, and registered under, U.S. securities laws and subject to oversight by the SEC.  In fact, before the stablecoin can be offered to investors, the SEC must clear the registration.  Once the stablecoin fund is fully registered, it will be available for sale for all investors in the U.S., regardless of accreditation status.  Further, since U.S. securities laws are generally considered the most stringent in the world, a registered U.S. offering would generally be a permissible investment for all retail investors in other jurisdictions outside the U.S. with minimal, if any, additional regulatory pre-conditions.


As an Investment Company Act fund, the fund will be required to provide periodic statements of performance and holdings to investors.  In addition, the fund will be subject to various investor protections such as the requirement to hold custody of assets with a qualified custodian, the requirement to engage a qualified transfer agent to keep track of shareholders, and the requirement to perform semi-annual and annual audits by a PCAOB auditor (i.e., an auditor approved by the SEC).  This assures investors that their stablecoins are fully-backed by investments held in institutions qualified to ensure their safety, and offers a level of transparency currently lacking in other stablecoin projects.


\subsection{D-SKRs and Class 8 Money Transmitter License on Lux MChain}

Lux MChain effectively utilizes a Class 8 Money Transmitter License to issue globally recognized, stable currencies backed by real-world assets, underpinned by Digital Safe Keeping Receipts (D-SKRs). This strategic application within the blockchain infrastructure enhances financial security and compliance.


\subsection{Class 8 Money Transmitter License: Key Facets for Lux MChain}

Diverse Financial Services: The license enables Lux MChain to provide a range of services, including bureau de change operations, payment services, cheque cashing, and importantly, the issuance of electronic money​​.


Electronic Money Issuance: Central to Lux MChain's capabilities, this license aspect allows for creating stable digital currencies backed by D-SKRs, aligning with global financial standards.


\subsection{D-SKRs in Asset-Backed Currency Issuance}

Asset Verification and Stability: D-SKRs serve as digital proof of ownership for underlying real-world assets, ensuring that the issued digital currencies are not only stable but also transparent and trustworthy.


Enhancing Digital Currency Reliability: The incorporation of D-SKRs on Lux MChain provides users with the assurance of asset backing, elevating the trust and reliability of the platform's digital currencies.


\subsubsection{Lux Dollar: A US Treasuries-Backed Stablecoin on Lux MChain}

\paragraph{Overview of Lux Dollar}

Lux Network introduces Lux Dollar, a novel stablecoin uniquely backed by short-term US Treasuries. This financial product is not just a stablecoin; it's also an investment vehicle, offering distinct advantages over current market offerings.


\subsection{Key Features of Lux Dollar}

Backed by US Treasuries: The backing by US Treasuries lends Lux Dollar unparalleled stability and security, a significant advancement from the typical structures of existing stablecoins.


Transformation of Stablecoin Landscape: With its actual yield and unmatched liquidity, Lux Dollar redefines the nature of stabletokens. It addresses the shortcomings of current stablecoins, which often face issues of being unregulated, expensive, and unredeemable.


Legal Structure and Practicality: Lux Dollar is meticulously designed with a strong legal structure, ensuring it operates effectively and compliantly in the financial space.


Unrivaled Liquidity and Accessibility: Projected to be the most liquid and universally accessible stable token, Lux Dollar aims to be free for consumers, setting a new standard in the stablecoin market.


\subsubsection{Muni Bonds Offering on Lux MChain}

\paragraph{Introduction to Muni Bonds Offering}

Lux's roadmap includes a unique offering in the realm of municipal bonds. This initiative is particularly tailored to meet the interests of the crypto community, showcasing the platform's capability to address specific investment needs and opportunities within this sector.


\paragraph{Highlighting Flexibility and Market Relevance}

Investment Opportunities in Municipal Bonds: Lux's offering in municipal bonds illustrates the platform's commitment to providing diverse and strategic investment options.


Adaptability to Market Needs: This move exemplifies Lux's adaptability, demonstrating its capacity to develop marketable and relevant financial products that resonate with the evolving preferences of the blockchain community.


\subsubsection{Digital Securities on Lux Platform}

\paragraph{Concept of Digital Securities}

Lux's initial product offerings are strategically designed to demonstrate the platform's advantages while catering to the nuanced requirements of the blockchain community.


\paragraph{Vision for Inclusive Financial Products}

Scalability Beyond Initial Products: Lux envisions transcending its initial offerings to embrace a broader spectrum of digital securities.


Accommodating a Wide Range of Investors: The goal is to launch, market, and distribute an array of products that align with the diverse return objectives and risk profiles of all investors, underlining Lux's commitment to inclusivity and versatility in investment solutions.


\section{Regulatory Compliance}

Lux is a multi-regulated broker under the stringent guidelines of several top-tier and high-end financial authorities. The Lux group has in its roadmap the ability to support clients from over 140 countries via several licensed Lux companies.


Each entity that has a license from a financial authority is required to comply with specific regulatory guidelines and these rules may vary from one regulator to another.


And because every company of Lux needs to conform with the applicable local regulatory framework, there could be differences in the available services and products that they offer for their clients based on their country of jurisdiction and the regulated entity that operates the platform. Be sure to read this page and our Lux analysis so you will be fully informed about this product.


\section{Current Licenses}

\subsection{North America}

Lux Card is issued by Metropolitan Commercial Bank (Member FDIC) pursuant to a license from Visa U.S.A. Inc. "Metropolitan Commercial Bank" and ``Metropolitan'' are registered trademarks of Metropolitan Commercial Bank.


Card Issuance


Exchange and Swap of assets


Crypto Exchange


Fractional Stock Investment


Lending


\subsection{UK}

CDAX, money transmission services level 8 in IOM covers digital money issuance and currency exchange.


CDAX License: https://www.iomfsa.im/media/2700/1389.pdf


Digital Money issuance


Exchange and Swap of assets


\subsection{Europe}

Blockbank has a crypto wallet, exchange and EMI license in Lithuania. Verne has been selected to help Lux acquire regulatory approvals and licenses internally.


Crypto Wallet and Exchange License


EMI License


\section{Research}

https://scalingbitcoin.org/milan2016/presentations/Scaling\%20Bitcoin\%203\%20-\%20Mark\%20Friedenbach.pdf


https://www.tik.ee.ethz.ch/file/a20a865ce40d40c8f942cf206a7cba96/Scalable\_Funding\_Of\_Blockchain\_Micropayment\_Networks\%20(1).pdf


\subsection{https://blockstream.com/bitcoin17-final41.pdf}

https://scalingbitcoin.org/milan2016/presentations/D1\%20-\%204\%20-\%20Andrew\%20Poelstra.pdf


\subsection{https://eprint.iacr.org/2016/1159.pdf}

\subsection{https://eprint.iacr.org/2018/104.pdf}

\subsection{https://github.com/ethereum/wiki/wiki/Sharding-FAQs}

\subsection{https://github.com/ethereum/wiki/wiki/Governance-compendium}

\subsection{https://tendermint.com/static/docs/tendermint.pdf}

\subsection{https://eprint.iacr.org/2017/634.pdf}

\subsection{https://cryptojedi.org/papers/dilithium-20170627.pdf}

\subsection{https://www.stellar.org/papers/stellar-consensus-protocol.pdf}

\subsection{https://bravenewcoin.com/assets/Whitepapers/QRL-whitepaper.pdf}

John R. Douceur, The Sybil Attack, International Workshop on Peer-to-Peer Systems (2002), pp. 251-260.


Satoshi Nakamoto, Bitcoin: A peer-to-peer electronic cash system (2008), https://bitcoin.org/ bitcoin.pdf


Leslie Lamport, Robert Shostak and Marshall Peaset, The Byzantine Generals Problem, ACM Transactions on Programming Languages and Systems Vol. 4, No. 3 (1982), pp. 382-401


Vitalik Buterin. Ethereum: A Next-Generation Smart Contract and Decentralized Application Platform (2013), https://github.com/ethereum/wiki/wiki/White-Paper


Delegated Proof-of-Stake Consensus, https://bitshares.org/technology/delegated-proof-of-stake-consensus/


Alan Demers, Dan Greene, Carl Hauser, Wes Irish, John Larson, Scott Shenker, Howard Sturgis, Dan Swinehart and Doug Terry, Epidemic algorithms for replicated database maintenance, PODC '87 Proceedings of the sixth annual ACM Symposium on Principles of distributed computing (1987), pp. 1-12.


Hyperledger Architecture, Vol. 1 (2017), https://www.hyperledger.org/wp-content/uploads/2017/08/Hyperledger Arch WG Paper 1 Consensus.pdf


Leemon Baird, The Swirlds Hashgraph Consensus Algorithm: Fair, Fast, Byzantine Fault Tolerance (2016), https://www.swirlds.com/downloads/SWIRLDS-TR-2016-01.pdf


David Mazirez, The Stellar Consensus Protocol: A Federated Model for Internet- level Consensus, Stellar Development Foundation (2016), https://www.stellar.org/papers/ stellar-consensus-protocol.pdf.


Gabriel Bracha and Sam Toueg, Asynchronous Consensus and Broadcast Protocols, Journal of the ACM, Vol. 32, Issue 4 (1985), pp. 824-840.


Miguel Castro and Barbara Liskov, Practical byzantine fault tolerance, Proceedings of the 3rd Symposium on Operating Systems Design and Implementation (1999), pp.173186.


Leslie Lamport, The Part-Time Parliament ACM Transactions on Computer Systems, Vol. 16, Issue 2 (1999), pp. 133-169.


Don Johnson, Alfred Menezes and Scott Van-Stone, The Elliptic Curve Digital Signature Algorithm (ECDSA), International Journal of Information Security, Vol. 1, Issue 1 (2001), pp. 36-63.


David L. Mills, Internet Time Synchronization: The Network Time Protocol, I.E.E.E. Trans. Comm. 39, No 10 (1991), pp. 1482-1493.


\subsection{https://blog.z.cash/snark-explain/}

\subsection{https://bitsofpy.blogspot.com/2014/03/homomorphic-encryption-using-rsa.html}

\subsection{https://bravenewcoin.com/assets/Whitepapers/QRL-whitepaper.pdf}

\end{document}
